% !TEX root = ../../tree_notes.tex
% =====================================================================
%  1.3  Kinds of Trees                                       sec:c01-kinds
%  Exports: def:c01-binary def:c01-full def:c01-perfect def:c01-complete
%           def:c01-balanced def:c01-skewed
%           thm:c01-perfect-count thm:c01-full-binary-leaves cor:c01-height-bounds
%           prob:c01-balanced prob:c01-count-nodes
% =====================================================================
\section{Kinds of Trees}\label{sec:c01-kinds}

\begin{previouslybox}
So far: the family vocabulary and the two theorems ($n-1$ edges, unique path), and one
DFS that pushes \emph{depth} in and returns \emph{height} and \emph{size} out
(\coderef{c01-measure}). This section asks the next question -- what happens when the
shape itself is \emph{promised} to you -- and every promise turns into a formula.

\begin{enumerate}[nosep, leftmargin=1.6em]
  \item With the edges convention, $\height(\nil)=$ \underline{\hspace{0.7cm}} and a lone
        node has height \underline{\hspace{0.7cm}}. \emph{(Ans: $-1$; $0$.)}
  \item In the running example of \figref{c01-depth-height}, what are $\height$ and
        $\size$ of the root $A$? \emph{(Ans: $\height(A)=2$, $\size(A)=6$.)}
  \item Which quantity is carried \emph{down} the recursion and which are returned
        \emph{up}? \emph{(Ans: depth down; height and size up.)}
\end{enumerate}
\end{previouslybox}

\begin{topicintro}
"Tree" hides a family of shapes, and each extra restriction buys a counting rule or a
faster algorithm. Knowing whether your tree is \emph{full}, \emph{complete}, \emph{perfect}
or \emph{balanced} tells you its height from its size, whether an array can store it, and
which problems become one line instead of a recursion.

\textbf{Prerequisites.} \secref{c01-anatomy} (height, size) and \thmref{c01-edges-nodes}.

\textbf{After this section you can}
\begin{itemize}[nosep]
  \item recognise and name binary, full, complete, perfect, balanced and skewed trees;
  \item never confuse "full", "complete" and "perfect" again;
  \item derive the node and leaf counts of perfect and full trees;
  \item bound the height of any binary tree with $n$ nodes by $\floor{\log_2 n}$ and $n-1$.
\end{itemize}
\end{topicintro}

\begin{recognitionbox}[Pattern card -- exploit the shape promise]
\reckey{``complete'', ``perfect'', ``full'', ``balanced'', ``height-balanced'', ``every
level filled except possibly the last'' -- any word that \emph{promises} a shape.}
\recuse{
  \item the problem guarantees a shape: use its counting rule instead of traversing
        (perfect subtree $=2^h-1$ nodes, \thmref{c01-perfect-count}; the $O(\log^2 n)$
        node count below is this idea at work);
  \item the problem \emph{asks} whether the shape holds (balanced, symmetric): a bottom-up
        check that returns per-node data and aborts early;
  \item a height bound lets you size arrays/stacks: $\floor{\log_2 n}\le h\le n-1$
        (\corref{c01-height-bounds}).
}
\recnotuse{
  \item no promise is stated -- never \emph{assume} completeness; a hidden skewed test
        case will break $O(\log n)$ logic;
  \item the task is to \emph{make} the tree balanced under updates -- that is a real
        self-balancing tree (\chref{c06});
  \item the shape is arbitrary but you need subtree aggregates -- tree DP (\chref{c10}).
}
One-line test: \emph{find the sentence that promises the shape, then ask what it makes
free: a formula, a bound, or a halving argument.}
\end{recognitionbox}

\subsection{Intuition}

A \term{binary tree} restricts every node to at most two children, conventionally called
\code{left} and \code{right}. That single rule is what makes tree algorithms branch in
two and lets us store a tree in an array. The other shapes in this chapter are binary
trees with one more promise each.

\begin{analogybox}[Single-elimination tournament]
A \emph{perfect} binary tree is a knockout bracket: two players per match (the children),
the winners fill the next round (level) completely, and the number of matches is fixed by
the number of players. Cut the last round short and you get a \emph{complete} tree; let
someone get a bye and it is no longer complete. The bracket picture is why the counts come
out as powers of two.
\end{analogybox}

\subsection{Definitions and Notation}

\begin{definition}{Binary and $m$-ary tree}{c01-binary}
A \term{binary tree} is a rooted tree in which each node has at most two children, named
\code{left} and \code{right} (the two slots are distinguished, so a node with one child
keeps track of whether it is the left or the right one). An \term{$m$-ary tree} allows at
most $m$ children; a general tree is an $m$-ary tree with unbounded $m$.
\end{definition}

\noindent\textbf{In plain words.} Two seats per node, and the seats are \emph{named}:
``one left child'' and ``one right child'' are \emph{different} binary trees. This one
nuisance decides inorder's meaning and half the reconstruction puzzles of \chref{c03}.

\begin{definition}{Skewed (degenerate) tree}{c01-skewed}
A binary tree is \term{skewed} (or \term{degenerate}) when every internal node has exactly
one child, so the tree is really a linked list: all-left or all-right. A skewed tree on $n$
nodes has height $n-1$ -- the worst case for any recursive tree code.
\end{definition}

\noindent\textbf{In plain words.} A tree that forgot to branch -- it is a linked list in
disguise. Whenever your solution is ``$O(n)$ time but recursion depth $n$'', this shape
is the judge's test case that kills it.

\begin{definition}{Full (proper) binary tree}{c01-full}
A binary tree is \term{full} (or \term{proper}, or strict) when every internal node has
\emph{exactly two} children -- a node has $0$ or $2$ children, never $1$.
\end{definition}

\noindent\textbf{In plain words.} No loners: every node you meet has either no children
at all or exactly two. One is not on the menu.

\begin{definition}{Perfect binary tree}{c01-perfect}
A binary tree is \term{perfect} when every internal node has two children \emph{and} all
leaves are at the same depth. It is a completely filled bracket of height $h$.
\end{definition}

\noindent\textbf{In plain words.} Fully packed, all the way down -- level $\ell$ holds
exactly $2^{\ell}$ nodes with zero gaps. The pyramid is perfect.

\begin{definition}{Complete binary tree}{c01-complete}
A binary tree is \term{complete} when every level except possibly the last is full, and
the last level is filled left to right with no gaps. Every perfect tree is complete; the
converse needs the last level to be full.
\end{definition}

\noindent\textbf{In plain words.} Cinema filling: every row packed left-to-right with no
empty seat in front of an occupied one; only the last row may be half full. This is
exactly the shape an \emph{array} can store without wasting slots -- which is why heaps
(\chref{c07}) and segment trees promise it.

\begin{definition}{Height-balanced tree}{c01-balanced}
A binary tree is \term{(height-)balanced} when for every node the heights of its left and
right subtrees differ by at most $1$. Balancedness is what keeps the height logarithmic so
searches stay fast; it is the property Chapters 5--6 fight to maintain.
\end{definition}

\noindent\textbf{In plain words.} No node may lean too hard: at every fork, one side may
deepen past the other by at most \emph{one} level. That tiny local promise keeps the
global height near $\log_2 n$ -- the difference between instant searches and walking a
list.

\begin{figure}[H]
\centering
\renewcommand{\arraystretch}{1.1}
\begin{tabular}{@{}c@{\quad}c@{\quad}c@{\quad}c@{}}
\begin{forest} snaptree [1 [2 [4][5]] [3 [6][7]]] \end{forest} &
\begin{forest} snaptree [1 [2 [4][5]] [3 [6]]] \end{forest} &
\begin{forest} snaptree [1 [2 [4][5]] [3 [6 [8][9]][7]]] \end{forest} &
\begin{forest} snaptree [1 [2 [3 [4]]]] \end{forest} \\[-3ex]
{\scriptsize perfect} & {\scriptsize complete, not full} &
{\scriptsize full, not complete} & {\scriptsize skewed} \\
{\scriptsize (also complete \& full)} & {\scriptsize (3 has one child)} &
{\scriptsize (6 has two, uneven)} & {\scriptsize (a list)}
\end{tabular}
\caption{Four shapes, four promises. Read the boxes under each to see which definition the
tree satisfies and which it breaks.}\label{fig:c01-shapes}
\end{figure}

\vsbox{Full vs Complete vs Perfect}{Full / Perfect}{Complete}{full: every internal node
has two children; perfect adds "all leaves level"}{every level filled left-to-right, last
level may be short but has no gaps}

\subsection{Counting Nodes and Leaves}

Each extra restriction pins down the size. These identities are the ones interviewers
expect cold.

\begin{theorem}{Perfect Tree Counts}{c01-perfect-count}
A perfect binary tree of height $h$ (counted in edges) has $2^{h+1}-1$ nodes, $2^{h}$
leaves, and $2^{\ell}$ nodes on level $\ell$ (levels numbered from $0$).
\end{theorem}
\begin{proofbox}
Level $\ell$ doubles each step (each node has two children), so it holds $2^{\ell}$ nodes;
levels run $0..h$, giving $\sum_{\ell=0}^{h} 2^{\ell} = 2^{h+1}-1$ nodes by the geometric
series. The only nodes with no children are those on the last level $h$, so leaves $=2^{h}$.
\end{proofbox}
\begin{insightbox}
Invert it: a perfect tree with $n$ nodes has $h = \log_2(n+1)-1$, so $n+1$ is a power of
two. This is why "is it perfect?" can be answered by counting levels, never by pointer
chasing.
\end{insightbox}

\begin{theorem}{Leaves of a Full Binary Tree}{c01-full-binary-leaves}
In a full binary tree with $I$ internal nodes and $L$ leaves, $L = I + 1$.
\end{theorem}
\begin{proofbox}
Count edges two ways. Every internal node has exactly two children, contributing two
parent-to-child edges, so the tree has $2I$ edges. But any tree on $n = I + L$ nodes has
$n-1$ edges (\thmref{c01-edges-nodes}). Therefore $2I = (I+L)-1$, i.e.\ $L = I + 1$.
\end{proofbox}

\noindent\textbf{Quick feel.} Start with one node: one leaf, zero internal -- the leaves
lead by exactly one. Every split converts a leaf into an internal node that sprouts
\emph{two} new leaves: leaves $-1+2$, internal $+1$ -- the gap of one never changes. The
proof above counts the same phenomenon through edges.

\begin{corollary}{Height Bounds of a Binary Tree}{c01-height-bounds}
A binary tree with $n$ nodes has height (in edges)
$\floor{\log_2 n} \le \height \le n-1$. The lower bound is attained by a perfect (or
complete) tree, the upper bound by a skewed tree.
\end{corollary}
\begin{proofbox}
A tree of height $h$ holds at most $2^{h+1}-1$ nodes (a perfect tree), so $n \le
2^{h+1}-1 < 2^{h+1}$, giving $h > \log_2 n - 1$, i.e.\ $h \ge \floor{\log_2 n}$. On the
other hand each edge can increase the longest root-to-leaf path by at most one, and there
are only $n-1$ edges, so $h \le n-1$. Both extremes are the perfect and skewed trees
above.
\end{proofbox}
\begin{insightbox}
"Balanced means height $\bigO{\log n}$." Keeping $h$ near its lower bound is the entire
point of AVL and red--black trees (\chref{c06}): it turns $O(n)$ lists into $O(\log n)$
searches.
\end{insightbox}

\usedin{\chref{c06} (self-balancing BSTs), \chref{c07} (heaps in an array) and \chref{c13}
(segment trees)}{complete and perfect trees are exactly the ones an array stores with no
wasted slots, and balanced trees are the ones whose height stays $\bigO{\log n}$.}

\subsection{Worked Examples}

\begin{example}{Reading the gallery}{c01-kinds-ex}
In \figref{c01-shapes} the third tree has internal nodes $\{1,2,3,6\}$ ($I=4$) and leaves
$\{4,5,8,9,7\}$ ($L=5$): indeed $L = I+1$, matching \thmref{c01-full-binary-leaves}. The
first is perfect with $h=2$, so $2^{3}-1 = 7$ nodes and $2^2 = 4$ leaves -- count them to
confirm.
\end{example}

\begin{tracebox}
For $n=6$ nodes, \corref{c01-height-bounds} says $\floor{\log_2 6}=2 \le h \le 5$. A
complete tree of six nodes has height $2$ (near the best), the skewed tree has height $5$
(the worst) -- both inside the band.
\end{tracebox}

\subsection{Variations and Twists}

\begin{itemize}
  \item \textbf{Verify the promise.} Instead of using completeness, \emph{test} it:
        LC 958 (Check Completeness) is level-order with a "once you see a gap, only
        leaves allowed" flag. \hfill\lv{2}
  \item \textbf{Array-shaped trees.} A complete tree lives in an array at indices
        $1..n$ with children at $2i,2i+1$ -- the address arithmetic behind heaps
        (\chref{c07}) and segment trees (\chref{c13}). Every "pointer" problem here
        has an array twin; find it as practice. \hfill\lv{3}
  \item \textbf{Balanced is a spectrum.} Height-balanced (AVL, \lv{3} preview
        \chref{c06}) versus weight-balanced versus merely "the depth difference of any
        two leaves is at most $1$" -- read which definition the statement uses; they
        are not equivalent.
        \hfill\lv{4}
  \item \textbf{$m$-ary generalisation.} Perfect $m$-ary trees count as
        $\sum_{\ell} m^{\ell} = (m^{h+1}-1)/(m-1)$ -- the same geometric series with a
        different ratio; K-ary heaps (the same array arithmetic as \chref{c07}) and
        K-ary tries reuse it.
        \hfill\lv{3}
  \item \textbf{Disguise.} "Tournament bracket", "file-system block", "heap snapshot",
        "memory segment layout" -- all complete/perfect vocabulary. The recognition
        sentence to hunt for: ``every level except possibly the last is completely
        filled''.
\end{itemize}

\subsection{Common Pitfalls}

\begin{warningbox}[Confusing full / complete / perfect]
They are not synonyms and not even comparable in one direction: a complete tree need not
be full (a last-level parent can have one child), and a full tree need not be complete
(uneven depth). "Perfect" is the only one that is both. Sketch a two-level and a
three-level example of each before you assert a property.
\end{warningbox}

\begin{warningbox}[Balanced is per-node, not just root]
Checking the height difference at the root only is not enough -- \probref{c01-balanced}
has test cases where the root looks balanced but a deep subtree is not. The condition must
hold for \emph{every} node.
\end{warningbox}

\smallskip\noindent\textbf{The trap as code.}

\begin{lstlisting}[caption={\textcolor{warningred}{WRONG -- root-only balance test}}, basicstyle=\ttfamily\small, breaklines=true]
bool isBalanced(TreeNode* r) {
  if (!r) return true;
  return abs(height(r->left) - height(r->right)) <= 1;
}
\end{lstlisting}
\begin{lstlisting}[caption={\textcolor{successgreen}{RIGHT -- per-node check via the $-1$ sentinel (first solved problem)}}, basicstyle=\ttfamily\small, breaklines=true]
bool isBalanced(TreeNode* r) { return check(r) != -1; }
// check() tests |lh-rh|<=1 at EVERY node while
// returning the height it already computed
\end{lstlisting}

\subsection{Problem Lab (fully solved)}

\begin{attemptgate}
Both problems below are top-$\lv3$ interview staples. Do \textbf{Balanced} with the
single-pass sentinel from memory (if you wrote two recursive functions, you built the
$\lv1$ rung, not the $\lv2$ one), then attempt \textbf{LC 222} cold -- the question
"what does \emph{complete} let me skip?" is the entire solve.
\end{attemptgate}

\solvedproblem{c01-balanced}{Balanced Binary Tree}{LeetCode 110}
  {https://leetcode.com/problems/balanced-binary-tree/}{LC Easy L2}
  {bottom-up height, early exit, sentinel}

\begin{problemstatement}[Balanced Binary Tree]
Given a binary tree, decide whether it is height-balanced: for every node, the heights of
its two subtrees differ by at most one.
\end{problemstatement}

\solstep{Understanding the Problem}
The naive plan -- a full height computation at every node -- is $O(n^2)$ (height is
$O(n)$ and we call it $n$ times). We can do one bottom-up pass that returns each node's
height and simultaneously checks its balance, reusing every child's height instead of
recomputing it.

\solstep{Key Observations}
\begin{observation}
If a subtree is unbalanced the whole tree is unbalanced, so we can abort early by
propagating a special "unbalanced" signal upward.
\end{observation}
\begin{observation}
Returning $-1$ as that signal works because a real height is never negative (with the
null $=0$ nodes convention here).
\end{observation}

\solstep{The Approach Ladder}
\begin{center}
\renewcommand{\arraystretch}{1.3}
\begin{tabular}{@{}c p{0.235\linewidth} p{0.28\linewidth} l c@{}}
\toprule
\textbf{Level} & \textbf{Approach} & \textbf{Idea} & \textbf{Time} & \textbf{Verdict} \\
\midrule
$\lv1$ & Recompute height at every node & definition straight from the book & $O(n^2)$ & \tle \\
$\lv2$ & Single-pass bottom-up with $-1$ sentinel & reuse each child's height & $O(n)$ & \ac \\
\bottomrule
\end{tabular}
\end{center}
Rung 1 is correct and worth saying out loud in an interview -- then collapse it to rung~2
by noticing that \code{height} is recomputed on subtrees that were already measured.

\solstep{Solution}
A recursive \code{check} returns the height of a balanced subtree, or $-1$ if that subtree
is unbalanced. A node is balanced iff both children returned non-negative and their
heights differ by at most one.
\begin{figure}[H]
\centering
\begin{tikzpicture}[diagram]
  \node[tnode] (r) at (0,0) {3};
  \node[tnode] (l) at (-1.6,-1.3) {9};
  \node[tnode] (c) at (1.6,-1.3) {20};
  \node[tnode] (cl) at (0.7,-2.6) {15};
  \node[tleaf] (cr) at (2.5,-2.6) {7};
  \node[tleaf] (x) at (0.7,-3.9) {6};
  \draw[tedge] (r) -- (l); \draw[tedgehl] (r) -- (c);
  \draw[tedge] (c) -- (cl); \draw[tedge] (c) -- (cr); \draw[tedgehl] (cl) -- (x);
\end{tikzpicture}
\caption{Not balanced: at node $20$ the left subtree (height $2$ through the orange edges
down to $6$) and the right subtree (height $0$) differ by more than one, even though the
root itself looks fine.}\label{fig:c01-balanced}
\end{figure}

\solstep{Proof of Correctness}
\begin{proofbox}[Induction on $\size(v)$]
By \thmref{c01-height-recursion} the returned height is correct whenever the subtree is
balanced. If a child returns $-1$, its subtree violates the condition at some node
(induction hypothesis) and $v$ correctly propagates $-1$. Otherwise $v$'s subtree is
balanced exactly when both children's subtrees are (induction) \emph{and} $v$ itself
satisfies $|\height(\text{left})-\height(\text{right})|\le 1$, which the code tests. The
root's return is negative iff the whole tree is unbalanced.
\end{proofbox}

\solstep{Implementation}
\begin{lstlisting}[caption={Balanced Binary Tree (LeetCode 110)}, label={lst:c01-balanced}]
#include <bits/stdc++.h>
using namespace std;

struct TreeNode {
    int val;
    TreeNode *left, *right;
    TreeNode(int x = 0) : val(x), left(nullptr), right(nullptr) {}
};

class Solution {
public:
    int check(TreeNode* r) {                 // height if balanced, -1 if not
        if (r == nullptr) return 0;
        int lh = check(r->left);   if (lh == -1) return -1;
        int rh = check(r->right);  if (rh == -1) return -1;
        if (abs(lh - rh) > 1) return -1;
        return 1 + max(lh, rh);
    }
    bool isBalanced(TreeNode* root) { return check(root) != -1; }
};
\end{lstlisting}

\solstep{Trace}
\begin{tracebox}
On \figref{c01-balanced}: \code{check(6)} returns $0$, so \code{check(15)}$=\max(0,\text{null }0)+1=1$;
\code{check(7)}$=0$; at $20$, $|1-0|=1$ ok, returns $2$; at $9$, both null, returns $0$;
at root $3$, $|0-2|=2>1$ returns $-1$. The root return is negative, so \code{isBalanced}
is \code{false}.
\end{tracebox}

\solstep{Complexity}
\begin{complexitybox}
Time $\bigO{n}$: each node is visited once and its children's heights are reused, not
recomputed (versus $O(n^2)$ for the root-only-height approach). Memory $\bigO{h}$ stack,
worst case $\bigO{n}$.
\end{complexitybox}

\solstep{Pitfalls and Variants}
\begin{warningbox}[Two traversals double the cost]
Writing \code{isBalanced} to call a separate \code{height} on every node is correct but
$O(n^2)$. The single-pass sentinel version above is what interviewers want to see.
\end{warningbox}
Variant: \emph{Minimum Absolute Difference in BST} or any per-node property uses the same
"return a value, abort early" bottom-up shape.

\begin{insightbox}
Key takeaway: fold two questions into one DFS -- \emph{return} the height and use an
impossible value ($-1$) as the "already broken" signal. The sentinel early-abort is the
reusable pattern for every per-node promise.
\end{insightbox}

\solvedproblem{c01-count-nodes}{Count Complete Tree Nodes}{LeetCode 222}
  {https://leetcode.com/problems/count-complete-tree-nodes/}{LC Medium L3}
  {complete-tree structure, perfect shortcut, divide and conquer}

\begin{problemstatement}[Count Complete Tree Nodes]
Given the root of a \emph{complete} binary tree, return the number of nodes. Aim for
better than $O(n)$: the structure should let you skip whole perfect subtrees.
\end{problemstatement}

\solstep{Understanding the Problem}
A plain DFS counting nodes is $O(n)$ and ignores the promise. In a complete tree, compare
the length of the left spine to the right spine of a subtree: if they are equal, that
subtree is \emph{perfect} and holds $2^{\text{height}}-1$ nodes
(\thmref{c01-perfect-count}) -- count it instantly. Otherwise recurse, but only down the
side that actually changes.

\solstep{Key Observations}
\begin{observation}
If the left-most and right-most paths from a node have the same length, the subtree is
perfect and its size is a power of two minus one.
\end{observation}
\begin{observation}
If they differ, exactly one child subtree is perfect, so each step halves the problem.
\end{observation}

\solstep{Think It Through With Me}
\begin{thinkbox}
\deadend{plain DFS over all $n$ nodes}{accepted -- but the statement \emph{promised}
completeness and you paid nothing for it; the judge's follow-up question is "why is your
answer $O(n)$?"}
\deadend{counting the last level by BFS-ing to depth $h$}{the full tree above the last
level has $2^h-1 = O(n)$ nodes -- visiting them to "find" the last level is the whole
$O(n)$ you promised to avoid.}
\aha{some \emph{subtrees} inside a complete tree are \emph{perfect} -- and a perfect
subtree's size is a closed formula. Detect one by comparing the left and right spines,
then never walk inside it again.}
\end{thinkbox}

\solstep{The Approach Ladder}
\begin{center}
\renewcommand{\arraystretch}{1.3}
\begin{tabular}{@{}c p{0.235\linewidth} p{0.28\linewidth} l c@{}}
\toprule
\textbf{Level} & \textbf{Approach} & \textbf{Idea} & \textbf{Time} & \textbf{Verdict} \\
\midrule
$\lv2$ & DFS counting every node & works, but ignores the shape promise & $O(n)$ & \ac \\
$\lv3$ & Perfect-subtree shortcut & spine test, $2^h-1$ count, halve each step & $O(\log^2 n)$ & \ac \\
\bottomrule
\end{tabular}
\end{center}
Rung 1 is accepted by the judge yet defeats the purpose of the problem -- rung~2 is the
intended solution and is the one to present.

\solstep{Solution}
Compare the left and right spine heights of the current subtree. Equal $\Rightarrow$
perfect $\Rightarrow$ return $2^{\text{spine}}-1$; otherwise return $1$ plus the counts of
the two (still-complete) children.
\begin{lstlisting}[caption={Count Complete Tree Nodes (LeetCode 222)}, label={lst:c01-count-nodes}]
#include <bits/stdc++.h>
using namespace std;

struct TreeNode {
    int val;
    TreeNode *left, *right;
    TreeNode(int x = 0) : val(x), left(nullptr), right(nullptr) {}
};

class Solution {
public:
    int leftHeight(TreeNode* n)  { int h = 0; while (n) { ++h; n = n->left;  } return h; }
    int rightHeight(TreeNode* n) { int h = 0; while (n) { ++h; n = n->right; } return h; }

    int countNodes(TreeNode* root) {
        if (root == nullptr) return 0;
        int lh = leftHeight(root), rh = rightHeight(root);
        if (lh == rh) return (1 << lh) - 1;              // perfect subtree: 2^lh - 1
        return 1 + countNodes(root->left) + countNodes(root->right);
    }
};
\end{lstlisting}

\solstep{Proof of Correctness}
\begin{proofbox}
When the two spine lengths agree ($=h$) the complete tree of height $h$ must have its last
level full, hence perfect, so the count is $2^{h}-1$ by \thmref{c01-perfect-count}. When
they differ, the root contributes $1$ and we recurse on the (still complete) children;
correctness follows by induction on height, since every subtree of a complete tree rooted
at a node is again complete.
\end{proofbox}

\solstep{Trace}
\begin{tracebox}
A perfect tree of height $2$ (7 nodes): at the root $lh=rh=3$, so we return $2^3-1=7$
without recursing. A complete tree of $6$ nodes: at the root $lh=3\ne rh$ (right spine
shorter), so we return $1+\text{count}(\text{left})+\text{count}(\text{right})$; the left
child is perfect ($2^2-1=3$) and the right has $2$, giving $1+3+2=6$.
\end{tracebox}

\solstep{Complexity}
\begin{complexitybox}
Spine lengths cost $O(\log n)$ each; we do this on one root-to-leaf path, so $O(\log^2 n)$
time -- a large win over $O(n)$ when the tree is deep and full. Memory $O(\log n)$ stack.
\end{complexitybox}

\solstep{Verification and Edge Cases}
\begin{itemize}[nosep]
  \item empty tree: returns $0$. \hfill\lv{3}
  \item perfect tree: shortcut fires at the root -- $O(\log n)$ total.
  \item $1 \ll lh$ with $lh=31$: overflows \code{int}; use \code{1LL << lh} and check the
        node-count constraint first (see the hardeningbox).
\end{itemize}

\solstep{Pitfalls and Variants}
\begin{warningbox}[$1 \ll lh$ overflow]
For very tall trees $2^{lh}$ can exceed \code{int}; if $n$ may be large use \code{long
long} and $1\ll lh$ only up to $lh<31$. Contest node counts keep $lh$ small, but check the
constraint.
\end{warningbox}
Variant: the same "measure both spines, shortcut perfect subtrees" idea finds the last
(level-order) node or the insertion point in a heap-backed array.

\begin{insightbox}
Key takeaways: (1) a promise is a coupon -- \emph{complete} buys the perfect-subtree
formula $2^{h}-1$; (2) detect a perfect subtree in $O(\log n)$ by comparing the two
spines; (3) each non-shortcut step halves the tree, giving $O(\log^2 n)$. Never assume
completeness the statement did not give you.
\end{insightbox}

\begin{hardeningbox}
\textbf{$\lv4$ sharpen:} the $O(\log^2 n)$ above re-measures spines at every level. Walk
down a \emph{single} root-to-leaf path instead, deciding left-vs-right at each node by
one spine comparison -- $O(\log n)$, the textbook optimum for this problem. The move
("spend $O(\log n)$ tests to choose a side, then commit") is the same binary-search
spirit you will meet in \chref{c13}. \textbf{$\lv4$ stress:} $1 \ll lh$ with
lh~$=31$: switch to \code{1LL~<<~lh} and check the constraint table first.
\end{hardeningbox}

\subsection{Practice Problems}

\begin{stuckbox}
\begin{itemize}[nosep, leftmargin=1.4em]
  \item Write the promise as a formula first: perfect $\Rightarrow 2^{h+1}-1$ nodes;
        full $\Rightarrow L = I+1$; complete $\Rightarrow$ array storage (\secref{c01-kinds}).
  \item Balanced checks fail most often because they test only the root -- the condition
        is \emph{per node}; return height and a sentinel together.
  \item If your solution never uses the promised shape, ask: which subtree can I count,
        skip, or halve in one step instead?
  \item $1 \ll h$ overflows for $h \ge 31$ -- check the node-count constraint before
        using the perfect-subtree shortcut.
\end{itemize}
\end{stuckbox}

\begin{practiceset}
  \practice{Symmetric Tree}{LeetCode 101}
    {https://leetcode.com/problems/symmetric-tree/}{LC Easy L2}
    {compare a tree with its mirror: recurse on (left.left, right.right) and (left.right, right.left).}
  \practice{Check Completeness of a Binary Tree}{LeetCode 958}
    {https://leetcode.com/problems/check-completeness-of-a-binary-tree/}{LC Medium L2}
    {BFS with a flag: after the first node with a missing child, every remaining node must be a leaf.}
  \practice{Binary Tree Level Order Traversal}{LeetCode 102}
    {https://leetcode.com/problems/binary-tree-level-order-traversal/}{LC Medium L2}
    {BFS with a queue; process one whole level per iteration to fill the levels.}
  \practice{Convert Sorted Array to BST}{LeetCode 108}
    {https://leetcode.com/problems/convert-sorted-array-to-binary-search-tree/}{LC Easy L2}
    {pick the middle element as root, recurse on halves -- builds a balanced tree. See \chref{c05}.}
  \practice{Binary Tree Right Side View}{LeetCode 199}
    {https://leetcode.com/problems/binary-tree-right-side-view/}{LC Medium L3}
    {Stretch: the last node of each BFS level; or a DFS that visits right before left.}
  \practice{Complete Binary Tree Inserter}{LeetCode 919}
    {https://leetcode.com/problems/complete-binary-tree-inserter/}{LC Medium L3}
    {Stretch: keep a queue of parents that still have a free slot to insert left-to-right.}
\end{practiceset}

% ---------------- spaced-repetition card ----------------
\begin{recallbox}
\selftest{
  \item Perfect tree of height $h$ (edges): nodes = \underline{\hspace{1.3cm}},
        leaves = \underline{\hspace{1.3cm}}.
  \item Full tree with $I$ internal nodes has \underline{\hspace{0.9cm}} leaves --
        and the proof counts \underline{\hspace{1.2cm}} two ways.
  \item Which shape promise lets you count nodes in $O(\log^2 n)$, and what exactly do
        you compare to spot a shortcut subtree?
  \item Give a tree that is full but not complete, and one that is complete but not full.
  \item Write the balanced-check sentinel \code{check} from memory, base case first.
}
Balanced means a \emph{per-node} condition; the sentinel return $-1$ collapses
``height'' and ``is-it-valid'' into one $O(n)$ DFS -- the same trick you will re-use
in every early-exit bottom-up problem of the book.
\end{recallbox}
